\documentclass[12pt,a4paper]{article}
\usepackage[utf8]{inputenc}
\usepackage[T1]{fontenc}
\usepackage{geometry}
\geometry{a4paper, left=2.5cm, right=2.5cm, top=2.5cm, bottom=2.5cm}
\usepackage{booktabs, tabularx, array, longtable}
\usepackage{float}
\usepackage{graphicx}
\usepackage{amsmath}
\usepackage{xcolor}
\usepackage{fancyhdr}
\pagestyle{fancy}
\fancyhf{}
\rhead{PlateVision --- Gestion de Projet Livrable L2}
\lhead{UCAC-ICAM / ULC-ICAM}
\cfoot{\thepage}
\setlength{\headheight}{14pt}
\usepackage{hyperref}
\hypersetup{
  pdftitle   = {Rapport de Gestion de Projet --- PlateVision},
  pdfauthor  = {\'{E}quipe PlateVision --- UCAC-ICAM / ULC-ICAM},
  colorlinks = true,
  linkcolor  = blue!70!black,
  urlcolor   = blue!70!black,
}

\title{\textbf{Rapport de Gestion de Projet}\\[0.5em]
\large PlateVision --- UCAC-ICAM / ULC-ICAM\\[0.3em]
\normalsize Livrable L2 --- §6 Cahier des Charges}
\author{\'{E}quipe PlateVision}
\date{\today}

\begin{document}
\maketitle
\tableofcontents
\newpage

% ── Section 1 : Organisation de l'équipe ─────────────────────────────────────
\section{Organisation de l'\'{e}quipe}
\label{sec:gp:equipe}

\begin{table}[H]
\centering
\caption{Composition de l'\'{e}quipe PlateVision --- r\^{o}les et modules}
\label{tab:gp:equipe}
\begin{tabular}{llll}
\toprule
\textbf{Membre} & \textbf{R\^{o}le} & \textbf{Module} & \textbf{Comp\'{e}tences mobilis\'{e}es} \\
\midrule
NKE GOUETH
  & Chef de projet
  & A1 (Naive Bayes)
  & Python, sklearn, coordination \'{e}quipe \\
ELIE NJOCK
  & D\'{e}veloppeur Senior
  & A2 (YOLO + OCR)
  & PyTorch, YOLOv8, EasyOCR, vision par ordinateur \\
ALI YOUSSOUF
  & Data Scientist
  & B (Clustering)
  & sklearn, matplotlib, r\'{e}duction dimensionnelle \\
FOUDA BASILE
  & Algorithmes
  & C (MDP)
  & NumPy, th\'{e}orie de la d\'{e}cision, MDP from scratch \\
NAPANI MA\"{E}L
  & Int\'{e}grateur / Rapporteur
  & D + coordination
  & \LaTeX{}, Git, analyse \'{e}thique, r\'{e}daction \\
\bottomrule
\end{tabular}
\end{table}

\paragraph{Gouvernance de l'\'{e}quipe}
NKE GOUETH assure la coordination quotidienne (stand-up 15 min chaque
matin, d\'{e}bloquer les blocages inter-modules). Les d\'{e}cisions
techniques importantes sont prises en commun (voir §\ref{sec:gp:decisions}).
Git est utilis\'{e} comme unique source de v\'{e}rit\'{e} (branche
\texttt{dataset\_nkegoueth}).

% ── Section 2 : Diagramme de Gantt ───────────────────────────────────────────
\section{Planification --- Diagramme de Gantt}
\label{sec:gp:gantt}

\begin{table}[H]
\centering
\caption{Diagramme de Gantt PlateVision --- 5 journ\'{e}es}
\label{tab:gp:gantt}
\begin{tabular}{p{5cm}ccccc}
\toprule
\textbf{T\^{a}che} & \textbf{J1} & \textbf{J2} & \textbf{J3} & \textbf{J4} & \textbf{J5} \\
\midrule
\multicolumn{6}{l}{\textit{Transversal}} \\
\quad Setup Git + environnement  & \textbf{X} & & & & \\
\quad DSD (Livrable L1)          & \textbf{X} & & & & \\
\quad Module D (carnet de bord)  & \textbf{X} & \textbf{X} & \textbf{X} & \textbf{X} & \\
\quad Rapport technique (r\'{e}daction) & & & & \textbf{X} & \\
\quad Gestion projet (Livrable L2) & & & & \textbf{X} & \\
\midrule
\multicolumn{6}{l}{\textit{Module A}} \\
\quad Pr\'{e}traitement donn\'{e}es      & \textbf{X} & & & & \\
\quad Naive Bayes (A1)           & \textbf{X} & \textbf{X} & & & \\
\quad YOLO + OCR (A2)            & & \textbf{X} & & & \\
\quad Comparaison A1 vs A2       & & \textbf{X} & & & \\
\midrule
\multicolumn{6}{l}{\textit{Module B}} \\
\quad CNN embeddings             & & & \textbf{X} & & \\
\quad K-Means clustering         & & & \textbf{X} & & \\
\quad Visualisation (PCA, t-SNE) & & & \textbf{X} & & \\
\quad Robustesse (k-analysis)    & & & \textbf{X} & & \\
\midrule
\multicolumn{6}{l}{\textit{Module C}} \\
\quad D\'{e}finition \'{e}tats + actions & & & & \textbf{X} & \\
\quad Matrice transitions + R    & & & & \textbf{X} & \\
\quad Value Iteration (VI)       & & & & \textbf{X} & \\
\quad Policy Iteration (PI)      & & & & \textbf{X} & \\
\quad Comparaison VI vs PI       & & & & \textbf{X} & \\
\midrule
\multicolumn{6}{l}{\textit{Soutenance}} \\
\quad D\'{e}mo live pipeline A$\to$B$\to$C & & & & & \textbf{X} \\
\quad Pr\'{e}sentation orale      & & & & & \textbf{X} \\
\bottomrule
\end{tabular}
\end{table}

\paragraph{Note sur la densit\'{e} J4}
La journ\'{e}e~4 est la plus charg\'{e}e : elle concentre le Module~C complet
et la r\'{e}daction des livrables. Ceci est intentionnel : le Module~C
(MDP) d\'{e}pend des r\'{e}sultats du Module~B (clusters), qui ne sont
disponibles qu'en J3. Le pipeline s\'{e}quentiel A$\to$B$\to$C impose
ce calendrier.

% ── Section 3 : Diagramme PERT ───────────────────────────────────────────────
\section{Diagramme PERT}
\label{sec:gp:pert}

\paragraph{D\'{e}pendances entre t\^{a}ches}

\begin{table}[H]
\centering
\caption{D\'{e}pendances PERT --- chemin critique PlateVision}
\label{tab:gp:pert}
\begin{tabular}{p{4cm}p{4cm}p{5cm}}
\toprule
\textbf{T\^{a}che} & \textbf{D\'{e}pend de} & \textbf{Justification} \\
\midrule
DSD             & ---             & Premier livrable, ind\'{e}pendant \\
Pr\'{e}traitement & DSD           & Dataset s\'{e}lectionn\'{e} avant traitement \\
A1 (Naive Bayes) & Pr\'{e}traitement & Caract\`{e}res pr\'{e}-d\'{e}coup\'{e}s requis \\
A2 (YOLO+OCR)  & Pr\'{e}traitement & Images 640$\times$640 requises \\
Comparaison A1/A2 & A1, A2       & Les deux modules termin\'{e}s \\
B (CNN embed.) & Comparaison A1/A2 & Meilleur mod\`{e}le s\'{e}lectionn\'{e} \\
B (K-Means)    & B (CNN embed.) & Embeddings g\'{e}n\'{e}r\'{e}s \\
B (Visualisation) & B (K-Means) & Clusters attribu\'{e}s \\
C (\'{E}tats)  & B (K-Means)    & cluster\_id = dimension principale de $\mathcal{S}$ \\
C (Transitions) & C (\'{E}tats) & $\mathcal{S}$ d\'{e}fini avant $P$ \\
C (R\'{e}comp.) & C (\'{E}tats) & $\mathcal{S}$ et $\mathcal{A}$ d\'{e}finis \\
C (VI + PI)    & C (Transitions), C (R\'{e}comp.) & MDP complet requis \\
Rapport C      & C (VI + PI)    & R\'{e}sultats num\'{e}riques n\'{e}cessaires \\
Soutenance     & Rapport C, D   & Tous modules termin\'{e}s \\
\bottomrule
\end{tabular}
\end{table}

\paragraph{Chemin critique}
\[
\text{A2 (YOLO)} \to \text{B (CNN)} \to \text{B (K-Means)}
\to \text{C (VI+PI)} \to \text{Rapport} \to \text{Soutenance}
\]

Ce chemin d\'{e}termine la dur\'{e}e minimale du projet. Tout retard
sur A2 ou B se propage directement au Module~C et aux livrables.
Le Module~D est parall\`{e}le \`{a} tous les autres modules (carnet de
bord quotidien).

% ── Section 4 : Matrice des risques ──────────────────────────────────────────
\section{Matrice des risques}
\label{sec:gp:risques}

\begin{table}[H]
\centering
\caption{Matrice des risques projet PlateVision}
\label{tab:gp:risques}
\begin{tabular}{clccp{4.5cm}}
\toprule
\textbf{ID} & \textbf{Risque} & \textbf{Probab.} & \textbf{Impact} & \textbf{Mitigation} \\
\midrule
R1 & Dataset non adapt\'{e} CEMAC
   & \'{E}lev\'{e}e
   & Critique
   & Augmentation cibl\'{e}e (formats CEMAC) + DSD transparent \\
R2 & GPU indisponible (Colab quotas)
   & Moyenne
   & \'{E}lev\'{e}
   & Fallback CPU pour Naive Bayes ; YOLOv8n (taille minimale) \\
R3 & OCR mauvaise performance sur plaques camerounaises
   & Moyenne
   & \'{E}lev\'{e}
   & Test de plusieurs librairies (EasyOCR, Tesseract, PaddleOCR) \\
R4 & MDP ne converge pas ($\varepsilon$ trop petit)
   & Faible
   & Moyen
   & $\varepsilon$, $\gamma$ ajustables en CLI (\texttt{--epsilon}, \texttt{--gamma}) \\
R5 & Code non ex\'{e}cutable en soutenance
   & Faible
   & Critique
   & Tests locaux J4 ; \texttt{demo.py} avec pipeline complet A$\to$B$\to$C \\
R6 & Conflit Git (merges)
   & Moyenne
   & Faible
   & Branche unique \texttt{dataset\_nkegoueth} + PR avant merge \\
R7 & Manque de donn\'{e}es plaques falsifi\'{e}es
   & \'{E}lev\'{e}e
   & Moyen
   & Simulation par homographie + bruit dans les labels OCR \\
\bottomrule
\end{tabular}
\end{table}

\paragraph{Mesure de contr\^{o}le globale}
La r\`{e}gle appliqu\'{e}e tout au long du projet : \textbf{aucun module
n'est d\'{e}clar\'{e} ``termin\'{e}'' sans qu'un script de test puisse
\^{e}tre ex\'{e}cut\'{e} depuis la racine du projet} (\texttt{python
modules/X/test\_X.py}). Cette r\`{e}gle a \'{e}t\'{e} valid\'{e}e par
l'ensemble de l'\'{e}quipe en J1.

% ── Section 5 : Suivi des décisions techniques ────────────────────────────────
\section{Suivi des d\'{e}cisions techniques}
\label{sec:gp:decisions}

\begin{table}[H]
\centering
\caption{Journal des d\'{e}cisions techniques PlateVision}
\label{tab:gp:decisions}
\begin{tabular}{p{1cm}p{4.5cm}p{4cm}p{4cm}}
\toprule
\textbf{Jour} & \textbf{D\'{e}cision} & \textbf{Alternatives reject\'{e}es} & \textbf{Justification} \\
\midrule
J1 & YOLOv8n (taille nano)
   & YOLOv8m, YOLOv8l
   & Temps r\'{e}el sur CPU MINT ; $\leq 6$\,ms/image \\
J1 & EasyOCR pour l'OCR
   & Tesseract, PaddleOCR
   & Meilleure gestion polices alphanumériques mixtes ; GPU optionnel \\
J1 & Git Flow branche unique
   & Branches par module
   & \'{E}quipe de 5 ; r\'{e}duction conflits ; livraison conjointe \\
J2 & $\gamma = 0.95$ pour le MDP
   & $0.90$, $0.99$
   & Horizon annuel ($\sim$20 \'{e}ch\'{e}ances) align\'{e} cycle vignettes MINT \\
J3 & $k = 3$ clusters (K-Means)
   & $k = 4$, $k = 5$
   & Elbow + silhouette ($s = 0.72$ pour $k=3$) + coh\'{e}rence m\'{e}tier \\
J3 & Pruning fr\'{e}quence $< 1\,\%$
   & Pruning $< 5\,\%$
   & Conserver \'{e}tats rares mais op\'{e}rationnellement importants (fraude) \\
J4 & MobileNetV2 pour embeddings
   & ResNet50, VGG16
   & Compromis taille/performance ; vecteurs 512 dim suffisants pour $k=3$ \\
J4 & \'{E}thique par design (garde-fou MDP)
   & Filtrage post-hoc
   & $R(s_{\text{conforme}}, a_{\text{saisie}}) \ll 0$ int\'{e}gr\'{e} dans le MDP \\
\bottomrule
\end{tabular}
\end{table}

% ── Section 6 : Historique Git ───────────────────────────────────────────────
\section{Historique Git}
\label{sec:gp:git}

\paragraph{Branche de travail}
Toutes les contributions ont \'{e}t\'{e} r\'{e}alis\'{e}es sur la branche
\texttt{dataset\_nkegoueth}, cr\'{e}\'{e}e \`{a} partir de \texttt{main}
en d\'{e}but de projet. Les commits suivent la convention
\texttt{Conventional Commits} (\texttt{feat}, \texttt{docs}, \texttt{fix},
\texttt{refactor}) pour assurer une traçabilit\'{e} s\'{e}mantique.

\paragraph{Politique de merge}
Les modifications sont int\'{e}gr\'{e}es par \textit{pull request} vers
\texttt{main} apr\`{e}s revue par au moins un autre membre de l'\'{e}quipe.
Aucun push direct sur \texttt{main} n'est autoris\'{e}.

\begin{table}[H]
\centering
\caption{Commits principaux --- branche \texttt{dataset\_nkegoueth}}
\label{tab:gp:commits}
\begin{tabular}{p{2.5cm}lp{7cm}}
\toprule
\textbf{Hash (abr.)} & \textbf{Auteur} & \textbf{Message} \\
\midrule
\texttt{bb80ffb} & andreonana & {[MODULE-C] feat(rewards): R(s,a) sourc\'{e}e FCFA} \\
\texttt{0c9ed65} & andreonana & {Module C: implement Value Iteration and Policy Iteration from scratch} \\
\texttt{9b25a37} & andreonana & {Module C: implement Policy Iteration from scratch} \\
\texttt{76bd886} & andreonana & {Module C: add VI vs PI comparison} \\
\texttt{3b2b3ec} & andreonana & {[INTERFACE] feat(demo): d\'{e}monstrateur temps r\'{e}el pipeline A$\to$B$\to$C} \\
\bottomrule
\end{tabular}
\end{table}

\paragraph{Statistiques de contribution}
L'historique complet est consultable via \texttt{git log -{}-oneline} sur
la branche \texttt{dataset\_nkegoueth}. La distribution des commits par
module refl\`{e}te l'effort relatif :
Module~C ($\sim$40\,\% des commits), Module~B ($\sim$30\,\%),
Module~A ($\sim$20\,\%), transversal -- documentation ($\sim$10\,\%).

\end{document}
